\documentclass[letterpaper,12pt]{article}
\usepackage{graphicx}
\usepackage{pdflscape}
\usepackage{amsmath}
\usepackage{amsthm}
\usepackage{lscape}
\usepackage{enumerate}
\usepackage{float}
\usepackage{bm}
\usepackage{grffile}
\usepackage{relsize}
\usepackage{tablefootnote}
\usepackage{footnote}
\usepackage{sectsty}
\usepackage{xcolor}
\usepackage{caption}
\usepackage{color}
\usepackage{natbib}
\usepackage{multirow}
\usepackage{bbm}
\usepackage[T1]{fontenc}
\usepackage[multiple]{footmisc}
\usepackage{appendix}
\usepackage{times}
\usepackage{relsize}
\usepackage{caption}
\usepackage{newtxtext,newtxmath}
\usepackage{enumitem}
\usepackage{hyperref}

%%%%%%%%%%%%%%%%%%%%%%%%%%%%%%%%%%%%%%%%%%%%%%%%%%%%%%%%%%%%
%%%   Title Page Information
%%%%%%%%%%%%%%%%%%%%%%%%%%%%%%%%%%%%%%%%%%%%%%%%%%%%%%%%%%%%
\title{\vspace{-.5in} ECON712: Replication Exercise}
\author{Leandro Freylejer}
\date{Report due: April 7th, 2026, in class \\ Presentation: TBA}

%%%%%%%%%%%%%%%%%%%%%%%%%%%%%%%%%%%%%%%%%%%%%%%%%%%%%%%%%%%%
%%%   Double and Single Space Commands
%%%%%%%%%%%%%%%%%%%%%%%%%%%%%%%%%%%%%%%%%%%%%%%%%%%%%%%%%%%%
\newlength{\defbaselineskip}
\setlength{\defbaselineskip}{\baselineskip}
\newcommand{\setlinespacing}[1]{\setlength{\baselineskip}{#1 \defbaselineskip}}
\newcommand{\doublespacing}{\setlength{\baselineskip}{1.3 \defbaselineskip}}
\newcommand{\singlespacing}{\setlength{\baselineskip}{1\defbaselineskip}}
\renewcommand{\thesubsection}{\arabic{subsection}}

%%%%%%%%%%%%%%%%%%%%%%%%%%%%%%%%%%%%%%%%%%%%%%%%%%%%%%%%%%%%
%%%   Set Margins Here
%%%%%%%%%%%%%%%%%%%%%%%%%%%%%%%%%%%%%%%%%%%%%%%%%%%%%%%%%%%%
\textwidth=6.5in
\textheight=9in
\oddsidemargin=0.10in
\evensidemargin=0.10in
\topmargin=-0.5in

%%%%%%%%%%%%%%%%%%%%%%%%%%%%%%%%%%%%%%%%%%%%%%%%%%%%%%%%%%%%%
%%%   Actual Document Begins Here
%%%%%%%%%%%%%%%%%%%%%%%%%%%%%%%%%%%%%%%%%%%%%%%%%%%%%%%%%%%%%
\begin{document}
\pagenumbering{arabic}

\maketitle
\vspace{-.3in}

\noindent \textbf{Overview.} Gaining applied research skills is a central learning outcome of this course. To that end, you are required to fully replicate the empirical results of a published paper that uses publicly available data. The exercise is composed of two parts: a written report (10\% of the final grade) and an in-class presentation (7\% of the final grade).\\

\noindent \textbf{Submission.} Your final submission must include: (i) the written report, (ii) all code used to download, clean, and analyse the data, and (iii) clear instructions for obtaining it, so that your results can be fully replicated. Please submit \textbf{annotated} do-files and log files.  The report should be between 8 and 12 pages, with 12pt font and 1.5 space between sentences (excluding code, tables, figures, and references).\\


%%%%%%%%%%%%%%%%%%%%%%%%%%%%
\section*{Part I: Written Report}
%%%%%%%%%%%%%%%%%%%%%%%%%%%%

The written report is structured around four components, described below. While you are expected to use your own words throughout, you should support your discussion with evidence from the paper and, where appropriate, cite additional references.

%--------------------------------------
\subsection*{1. \quad Summary and Contribution (15 points)}
%--------------------------------------

In this section, provide a concise but complete summary of the paper. Your discussion should address the following questions:

\begin{enumerate}[label=(\alph*)]
    \item \textit{Research question.} What is the central question the paper seeks to answer? What is the economic or policy motivation for studying this question?

    \item \textit{Contribution to the literature.} How does the paper relate to the existing literature on this topic? What gap does it fill, and what advantages does its approach have over prior work? Are there any limitations relative to other approaches you are aware of?

    \item \textit{Main findings.} What are the paper's main empirical results and conclusions? Do the authors claim that their estimates are causal? If so, on what basis?
\end{enumerate}

%--------------------------------------
\subsection*{2. \quad Identification Strategy (20 points)}
%--------------------------------------

This section is the core of the written report. You should provide a detailed explanation of how the authors identify the causal effect of interest. Your discussion should address the following:

\begin{enumerate}[label=(\alph*)]
    \item \textit{The identification problem.} What is the fundamental identification challenge the paper faces? What sources of endogeneity or selection bias would arise from a naive comparison?

    \item \textit{The identification strategy.} Explain clearly the identification strategy used by the authors (e.g., randomized experiment, instrumental variables, difference-in-differences, regression discontinuity, matching). What assumptions are required for the strategy to deliver causal estimates? Are these assumptions plausible in the context of the paper? Discuss any evidence the authors provide in support of these assumptions.

    \item \textit{Empirical specification.} Present the main regression specification(s) of the paper that you are replicating. Define all variables and parameters. Explain why the authors include the controls they do, and discuss whether any important variables may have been omitted.

    \item \textit{Critical assessment.} In your view, does the identification strategy convincingly establish causality? Are there any threats to identification that the authors do not fully address? Be specific.
\end{enumerate}

%--------------------------------------
\subsection*{3. \quad Replication (40 points)}
%--------------------------------------

In this section, you will replicate the main empirical results of the paper. Specifically:

\begin{enumerate}[label=(\alph*)]
    \item \textit{Data.} Describe the dataset(s) used in the paper. Where is the data publicly available? Explain how you obtained, merged, and cleaned the data. Discuss any challenges you encountered in this process and how you resolved them. Attach all data cleaning and analysis code to your submission.

    \item \textit{Descriptive statistics.} Reproduce the main descriptive statistics table(s) from the paper. Comment briefly on the key features of the data.

    \item \textit{Main results.} Replicate the paper's main regression table(s) and figure(s). Are your results consistent with those reported in the paper? If there are discrepancies, discuss possible explanations.

\end{enumerate}

\noindent \textbf{Grading note:} Minor discrepancies with the results in the paper due to data version differences or rounding will not be penalized, provided you explain them clearly.

%--------------------------------------
\subsection*{4. \quad Critical Discussion and Proposed Extensions (25 points)}
%--------------------------------------

The final section of the report is reserved for your own critical assessment of the paper and your ideas for improving or extending the analysis.

\begin{enumerate}[label=(\alph*)]
    \item \textit{Assessment of empirical choices.} Discuss the authors' main empirical choices: the choice of outcome variable, sample restrictions, control variables, standard error clustering, and the form of the regression (e.g., linear, log-linear). Do you agree with these choices? Are there alternatives that would have been more appropriate or informative?

    \item \textit{Proposed improvements.} Describe at least two concrete ways in which you would improve or extend the analysis if you were to write a follow-up paper. These could include: using a different or richer dataset, applying a different identification strategy, addressing a heterogeneity in treatment effects, or studying a related but unexplored research question. For each proposed improvement, explain (i) what additional insight it would provide relative to the original paper, and (ii) what data or methods would be required to implement it.

    \item \textit{Overall assessment.} Do you find the paper's conclusions convincing? Would you have approached the research question differently? Is there anything you particularly liked or disliked about the analysis?
\end{enumerate}

%%%%%%%%%%%%%%%%%%%%%%%%%%%%
\section*{Part II: Presentation}
%%%%%%%%%%%%%%%%%%%%%%%%%%%%

On TBA, each student will present their replication exercise to the class. Presentations should be around 15 minutes, followed by 5 minutes of questions. Your presentation should follow the same structure as the written report:

\begin{enumerate}
    \item A brief summary of the paper and its contribution (3--4 minutes).
    \item A clear explanation of the identification strategy (3-- minutes).
    \item A discussion of the main replicated results, with at least one table or figure (4--5 minutes).
    \item Your proposed improvements and overall assessment (3--4 minutes).
\end{enumerate}

\noindent You are required to use slides. Make sure your tables and figures are clearly readable. You will be evaluated on the clarity of your presentation, the depth of your understanding of the paper, and the quality of your critical assessment.

%%%%%%%%%%%%%%%%%%%%%%%%%%%%
\section*{List of Papers}
%%%%%%%%%%%%%%%%%%%%%%%%%%%%

The following is a list of assigned papers.

\begin{itemize}

\item \textit{Abel Ocheng:} Pretis, Felix (2022) ``Does a carbon tax reduce CO2 emissions? Evidence from British Columbia.", \textit{Environmental and Resource Economics}, 83(1), 115-144. \\
Topic: Carbon Tax and emissions. Replication Exercise: Replicate all tables and graphs up to page 131 (\textbf{do not include} synthetic control and robustness checks in appendices).

\item \textit{Ayomide Oladini:} Krueger, A.B. (1999). ``Experimental Estimates of Education Production Functions.'' \textit{Quarterly Journal of Economics}, 114(2), 497--532.\\
    \textit{Topic: Returns to class size. Replication Exercise: Replicate all tables and graphs up to page 519 (\textbf{do not include} Section C onwards) . Data:  Project STAR available @ Harvard Dataverse.}

\end{itemize}

\noindent If you have further questions about the assignment or the paper selection, please do not hesitate to contact me.

\end{document}
